\anexo{Guiones de Entrevistas}

\section{Guiones de Entrevistas}

\subsection{Entrevista 1 -- Agente Inmobiliaria}

\textbf{Encuestador:} Hola, Sharon. Muchas gracias por tu tiempo. Somos estudiantes de Ingeniería Informática y estamos haciendo una investigación sobre el funcionamiento del mercado inmobiliario en CABA. Nos interesa conocer tu experiencia real del día a día, así que sentite libre de explayarte. ¿Te molesta si grabamos el audio solo para fines académicos y de transcripción interna?

\textbf{Entrevistada:} No, para nada, adelante.

\textbf{Encuestador:} Perfecto, te lo agradecemos. Para arrancar, ¿hace cuánto estás en el rubro y cómo es un día típico de trabajo para vos?

\textbf{Entrevistada:} Estoy en el rubro hace 3 años ya y cada día aprendo algo nuevo. Un día típico es entrar a la oficina, preparar y acomodar todo para que esté presentable, así a la hora de tener firmas o recibir clientes que vienen a pagar los alquileres esté todo prolijo y ordenado. Después salgo a la calle. Si tengo alguna visita voy directo para allá, y sino voy entrando a locales con tarjetas de nuestra empresa y voy entregándolas. A la hora de estar con clientes tratamos de tener información previa antes de la visita o el encuentro.

\textbf{Encuestador:} Interesante. Y cuando les llega una propiedad nueva, ¿cómo es ese proceso?

\textbf{Entrevistada:} Es bastante más simple de lo que se cree. Nos llega la propiedad, le preguntamos al propietario qué precio quiere poner y cuál es el valor mínimo por el que está dispuesto a rebajarlo. Charlamos un poco, y si vemos que está fuera de valor y el dueño no está dispuesto a bajar, directamente no aceptamos la propiedad. Si está dispuesto, llegamos a un acuerdo en un precio razonable y la publicamos. Para eso usamos datos de otras propiedades de la misma zona con características similares, como los metros cuadrados, la cantidad de ambientes, los amenities, y también la experiencia que uno va acumulando. Si el propietario no cede, no publicamos ni trabajamos esa propiedad.

\textbf{Encuestador:} ¿Y qué tipo de información te suelen pedir los clientes cuando se interesan por una propiedad?

\textbf{Entrevistada:} Siempre preguntan a cuánto está dispuesto a vender o alquilar realmente el dueño, o si pueden hacer una oferta con lo que ellos consideran que vale. Les sirve mucho saber en qué rondan los precios en esa zona para características similares, si nosotros creemos que está fuera de valor, cuántas veces se mostró la propiedad o hace cuánto está publicada.

\textbf{Encuestador:} Claro, la zona también debe influir bastante. ¿Qué tanto peso le dan al barrio o al entorno?

\textbf{Entrevistada:} Juega un papel muy importante. Si es una zona con mucha oferta y poca demanda, por ejemplo, eso nos condiciona bastante. También hay zonas más caras que otras, por eso el entorno siempre es clave. Lo primero que preguntan los clientes es si hay transporte público de fácil acceso o si es una zona transitada. Después si hay supermercados, hospitales o escuelas cerca, y por último si tiene plazas o shoppings en los alrededores. Nosotros tenemos la oficina en Palermo y todos conocemos bien el barrio, así que es fácil responder esas consultas. Cuando nos llega una propiedad en otro barrio, tratamos de ir a recorrer la zona antes de mostrarla, o le pedimos información al propietario.

\textbf{Encuestador:} ¿Y cómo hacen para transmitirle al cliente cómo es el entorno de la propiedad?

\textbf{Entrevistada:} Tratamos de mostrarles las calles principales que la rodean. En Palermo, por ejemplo, siempre mencionamos la Avenida del Libertador, Alcorta, las plazas cercanas, algún bar conocido de la zona. Siento que me vendría bien poder hacer un pequeño tour por los alrededores para que el cliente lo vea de verdad. Algo que me encantaría incorporar es mostrárselo directamente por Google Maps, que es algo que hoy no hago pero que sumaría mucho.

\textbf{Encuestador:} Tiene mucho sentido. Cambiando un poco el tema, ¿qué lugar ocupan las aplicaciones en tu trabajo diario?

\textbf{Entrevistada:} Son herramientas esenciales, y creo que más todavía en este rubro, porque nos permiten ver qué hay disponible en el mercado de forma actualizada, publicar propiedades y comparar precios o tasaciones en cualquier momento, desde la oficina o desde casa. Lo que más me resulta útil es poder consultar y comparar cuando quiero, sin restricciones. Y creo que hay mucho por mejorar: siento que una aplicación podría, según las características del inmueble, usar inteligencia artificial para determinar el valor real de la propiedad y hacer la tasación por nosotros.

\textbf{Encuestador:} Para cerrar, ¿cuáles dirías que son los principales desafíos del sector hoy?

\textbf{Entrevistada:} El mayor desafío, para mí, es la cantidad de inmobiliarias que existe, y que cualquier persona, tenga matrícula o no, puede publicar un inmueble en cualquiera de las plataformas. Eso complica bastante el panorama. Aunque lo veo como un obstáculo que, honestamente, no tiene solución fácil.

\textbf{Encuestador:} Buenísimo, Sharon, muchísimas gracias. Nos fue de gran ayuda todo lo que nos contaste.

\subsection{Entrevista 2 -- Propietaria de varios inmuebles en CABA}

\textbf{Encuestador:} Para empezar y ponernos en contexto, contanos un poco cómo está compuesto tu portfolio actualmente y cómo te organizás en el día a día para gestionar los alquileres.

\textbf{Entrevistada:} Actualmente tengo cuatro departamentos: un monoambiente, dos de 2 ambientes y uno de 3 ambientes. La gestión la hacemos en colaboración con mi familia. A principio de mes enviamos la facturación detallada que incluye el alquiler y los gastos fijos, ya que nos ha pasado que los inquilinos se olvidan de pagar servicios básicos y luego tenemos que asumir nosotros los inconvenientes. Encuestador: Mencionaste que algunas propiedades las manejás de forma directa y otras con inmobiliaria. ¿Cómo decidís qué delegar y qué manejar por tu cuenta?

\textbf{Entrevistada:} Priorizamos siempre al inquilino directo, principalmente por recomendaciones y confianza, y porque resulta más económico para el inquilino. Si no logramos alquilarlo por dueño directo ---donde tiene que haber un cierto nivel de confianza---, recurrimos a una inmobiliaria para que asuma ese riesgo.

\textbf{Encuestador:} ¿Cómo es tu experiencia general trabajando con inmobiliarias para la administración o publicación de tus inmuebles?

\textbf{Entrevistada:} Para la publicación son muy rápidas; cuando trabajamos con ellos los inmuebles no pasan mucho tiempo vacíos. Sin embargo, en lo que respecta a la administración, se desentienden (``se lavan las manos''). Hacen el nexo para firmar el contrato y presentar a las partes, pero si luego hay algún percance, la inmobiliaria no se hace responsable y tenés que arreglarlo directamente con el inquilino.

\textbf{Encuestador:} Cuando una propiedad tuya se publica en portales tradicionales (como Zonaprop o Argenprop), ¿qué uso le das vos a esas plataformas y qué información sentís que falta?

\textbf{Entrevistada:} Aún no hemos publicado como dueños directos allí. Lo evaluamos como alternativa, pero nos da temor la falta de confianza sobre el inquilino (si va a cuidar el inmueble o si va a pagar en tiempo y forma). Estos portales son buenos para presentar inmuebles, pero tienen una gran debilidad: no brindan información para validar la confiabilidad o seguridad del inquilino.

\textbf{Encuestador:} A la hora de renovar un contrato o buscar un nuevo inquilino, ¿cómo es tu proceso para definir el valor que vas a pedir y de dónde sacás las referencias?

\textbf{Entrevistada:} Nos fijamos en las plataformas inmobiliarias, miramos los precios de la zona y las características de inmuebles similares al nuestro para tener un valor aproximado. Luego coordinamos con el inquilino los aumentos (cada 3 o 6 meses). Somos bastante flexibles; si tenemos un buen inquilino y buenas referencias, preferimos ganar un poco menos pero asegurarnos de que el inmueble se mantenga en buen estado.

\textbf{Encuestador:} Nos comentaron que estás considerando vender una de tus propiedades. ¿Cómo fue el proceso de tasar ese inmueble para la venta y cómo validaste ese número?

\textbf{Entrevistada:} Pedimos una tasación a una inmobiliaria tradicional para un departamento en el Microcentro. Validamos que el precio era correcto porque nos entregaron un informe muy detallado que evaluaba metros cuadrados, iluminación y, fundamentalmente, un registro histórico de ventas de otros años con el que podíamos evaluar la curva de precios. Igual, todavía no lo vendimos porque estamos evaluando si es el momento oportuno. Una cosa es el precio propuesto y otra encontrar a alguien dispuesto a pagarlo en el mercado actual.

\textbf{Encuestador:} ¿Te ha pasado de tener discrepancias de precio en tus propiedades?

\textbf{Entrevistada:} En ventas no, pero en alquileres me ha pasado de firmar contratos con porcentajes de aumento predefinidos que, con el tiempo, se desfasaron demasiado y el alquiler quedó a un precio muy bajo respecto al mercado. Al manejar los bienes de forma individual, a veces uno no accede tan fácil a la información actualizada que sí tienen las inmobiliarias.

\textbf{Encuestador:} El mercado tiene momentos donde conviene alquilar y otros vender. ¿Qué factores mirás para decidir sacar una propiedad del alquiler y ponerla a la venta?

\textbf{Entrevistada:} Observo mucho el auge y la popularidad de las zonas. Hay zonas que tienen un ``boom'' temporal y luego esa popularidad se traslada a otro barrio. Cuando noto que la popularidad de una zona recae, considero que es mejor vender e invertir ese capital en otro negocio o en una zona con una demanda más estable.

\textbf{Encuestador:} Si tuvieras el capital para comprar una nueva propiedad en CABA, ¿qué es lo primero que analizás de un barrio o propiedad para saber si es buen negocio?

\textbf{Entrevistada:} Primero la popularidad de la zona. Luego me fijo en los amenities del edificio (hacen que sea más fácil de alquilar) y en el tamaño (los departamentos grandes son muy difíciles de alquilar por los costos, la gente suele comprarlos). En cuanto al entorno, busco seguridad y transporte; el perfil que alquila en Capital (estudiantes o trabajadores) necesita movilidad fácil, vida nocturna y no quiere estar lejos de sus actividades.

\textbf{Encuestador:} ¿Cómo investigás esos datos si es un barrio en el que nunca viviste?

\textbf{Entrevistada:} Busco en la web, leo noticias y visito el barrio. Sin embargo, mi herramienta número uno es el ``boca a boca''. Es mucho más confiable hablar con conocidos, ya que la información en internet suele ser muy subjetiva o difícil de validar. Hoy no existe una página centralizada donde puedas ver todos los detalles reales de un barrio.

\textbf{Encuestador:} En base a tu experiencia, si pudieras tener una herramienta ideal para administrar tus inversiones y tomar decisiones, ¿qué problemas te gustaría que resuelva?

\textbf{Entrevistada:} Me gustaría que tuviera las siguientes funciones:

\begin{itemize}
\item \textbf{Validación de inquilinos:} Un sistema de referencias de propietarios anteriores (similar a Airbnb, pero para largo plazo) para asegurar la calidad y el cuidado del inmueble.
\item \textbf{Centralización de la gestión:} Tener toda la administración unificada en una sola plataforma para no tener que usar múltiples canales por fuera.
\item \textbf{Tasación automática (Rango de precios):} Que la herramienta me brinde un rango de precios predefinido por zona para saber exactamente cuánto vale mi departamento sin ``regalarlo ni aprovecharme''.
\item \textbf{Gestión legal integrada:} Que permita tener modelos de pre-contrato y la posibilidad de realizar firmas digitales para las renovaciones, evitando tener que reunirme presencialmente con los inquilinos.
\end{itemize}

\textbf{Encuestador:} Muchísimas gracias por tu tiempo, nos será de gran ayuda todo lo que aportaste.

\subsection{Entrevista 3 -- Persona en Proceso de Compraventa}

% TODO: contenido pendiente en el documento Word original.
